\documentclass[11pt, addpoints]{exam}
\usepackage{amsmath}
\usepackage{amsthm}
\usepackage{amsfonts}
\usepackage{amssymb}
\usepackage{mathrsfs}
\checkboxchar{$\Box$}
\usepackage{bbding}
\usepackage{pifont}
\usepackage{wasysym}
\usepackage[shortlabels]{enumitem}
\renewcommand{\qedsymbol}{$\blacksquare$}
\RequirePackage{amssymb, amsfonts, amsmath, latexsym, verbatim, xspace, setspace}
\usepackage{array} %table width
%\RequirePackage{tikz, pgflibraryplotmarks}
\usepackage[margin=1in]{geometry}
\usepackage{blindtext}
\usepackage{multicol}
\setlength{\columnsep}{1cm}
\usepackage{afterpage}
\usepackage{color} 

\newcommand\blankpage{%
    \null
    \thispagestyle{empty}%
    \addtocounter{page}{-1}%
    \newpage}
%% https://tex.stackexchange.com/a/203986
% Create a True-False question format
\newcommand*{\TrueFalse}[1]{%
\ifprintanswers
    \ifthenelse{\equal{#1}{T}}{%
        \textbf{TRUE}\hspace*{14pt}False
    }{
        True\hspace*{14pt}\textbf{FALSE}
    }
\else
    {True}\hspace*{20pt}False
\fi
} 
%% The following code is based on an answer by Gonzalo Medina
%% https://tex.stackexchange.com/a/13106/39194
\newlength\TFlengthA
\newlength\TFlengthB
\settowidth\TFlengthA{\hspace*{1.16in}}
\newcommand\TFQuestion[2]{%
    \setlength\TFlengthB{\linewidth}
    \addtolength\TFlengthB{-\TFlengthA}
    \parbox[t]{\TFlengthA}{\TrueFalse{#1}}\parbox[t]{\TFlengthB}{#2}}
%%

%%
\pagestyle{headandfoot}
\extraheadheight{.25in}
\firstpageheader{}
{\bfseries\large Indian Institute of Technology Mandi\\
End-Semester Exam\\
MA524 Probability \& Statistics\\
Instructor:\enspace Prof. Syed Abbas}{} %%%Main header
\runningheader{\bfseries\large MA524 Probability \& Statistics\\Prof. Syed Abbas}
{}
{ End-semester Exam\\May 17, 2024} %%running header
\firstpagefooter{}{Please go on to the next page\ldots}{Page \thepage\ of \numpages} %%%first page footer
\runningfooter{}{}{Page \thepage\ of \numpages} %%%running footer
%\runningfootrule
\begin{document}
\vspace{0.2in}
Roll Number\enspace \fillin \hfill Name:\enspace \hrulefill \\
\vspace{0.1in}
%\par Date: May 17, 2024\hfill Instructor:\enspace Prof. Syed Abbas
%\vspace{0.1in}

%%%%%%% Instructions begin here
\begin{center}
\fbox{\fbox{\parbox{5.5in}{\centering
\vspace{0.2in}
   \textbf{Please read the following instructions carefully.}
\begin{itemize}
    \item This question booklet contains \numquestions\ questions, \numpages\ pages (including the cover) for the total of \numpoints\ points. All questions are compulsory. You have 180 minutes to complete this exam.
    \item \textbf{Write your answers for Question No. \ref{que-1}, \ref{que-2}, and \ref{que-3} on question booklet}. Don't write any additional justification for these questions in the answer booklet.
    \item \textbf{Write your answers for Question No. \ref{que-4}, \ref{que-5} and \ref{que-6} on the answer booklet.} You can ask for additional sheets if required. Please write the full steps involved in your proof/solution; partial credit may be awarded for partial steps.
      
\end{itemize}
   }}}
\end{center}
%%%%% Grade table here
\begin{center} 
\gradetable[v][questions]
\end{center}
\hspace{0.5in}
\vspace{0.3in}
\begin{questions}
\addpoints
%\uplevel{Questions \ref{que-1} through~\ref{que-3} should be on answered on (this) question booklet.}
\newpage
\question[5]{Mark each of the following statement TRUE or FALSE.%}% ; \textbf{NO JUSTIFICATION IS REQUIRED.}
\begin{parts} 
   \part \TFQuestion{F}{Let $X$ be the number of heads in two tosses of a coin. If the coin is fair, the distribution of $X$ is then
  \begin{center}
      \begin{tabular}{|m{5em}||m{2em}|m{2em}|m{2em}|}
        \hline
          $x$ &  0 & 1 & 2\\ [0.5ex]
          \hline 
          $P\{X=x\}$ & $\frac{1}{4}$ & $\frac{2}{4}$ & $\frac{1}{2}$ \\ [0.5ex]
          \hline 
      \end{tabular}.
  \end{center}
   } % 4.1 (e) http://www.stat.ucla.edu/~dinov/courses_students.dir/02/Fall/STAT100A.dir/STAT_100.dir/Stat100FinalReview.pdf
   
   \part \TFQuestion{F}{Let $X$ be the number of heads in two tosses of a coin. If the coin is fair, the distribution function of $X$ is then
   \begin{equation*}
       F(x) =\begin{cases}
           0.25, & x<0 \\
           0.75 & 0 \leq x \leq 1 \\
           1 & 1 \leq x
       \end{cases}.
   \end{equation*}
   }
   \part \TFQuestion{T}{If $X$ is a binomial random variable, then $E[X]=np$.} 
   %http://www.stat.ucla.edu/~dinov/courses_students.dir/02/Fall/STAT100A.dir/STAT_100.dir/Stat100FinalReview.pdf
   \part \TFQuestion{T}{Let $X\sim U[49,52]$, the probability density function is given by the piecewise function,
   \begin{equation*}
       f(x)=\begin{cases}
           0 & \text{if } x <49 \\
           \frac{1}{3} & \text{if } 49 \leq x \leq 52 \\
           0 & \text{if } x> 52.
       \end{cases}
   \end{equation*}
   }
   \part \TFQuestion{T}{If $X \sim N(0,1)$ then $P\{X \leq x\}=P\{X <x\}$.}
\end{parts} 
} \label{que-1}
\question[5]{Multiple choice questions. There is a single correct answer to each of the following questions. \textbf{Select the correct answer (\CheckedBox)} % ; NO JUSTIFICATION IS REQUIRED}.
\begin{parts}
    \part Out of $6$ unbiased coins, $5$ are tossed independently and they all result in heads. If the $6^{th}$ is now independently tossed, the probability of getting head is? \\
    \begin{oneparcheckboxes}
        \choice $1$.
        \choice $0$.
        \correctchoice $\frac{1}{2}$.
        \choice $\frac{1}{6}$.
    \end{oneparcheckboxes}
    % CSIR NET 2018 (II) Dec Part A Que 2
    % https://csirhrdg.res.in/SiteContent/ManagedContent/ContentFiles/20190705142006547mathA_Dec2018.pdf
    %\vspace*{\stretch{1}}
    \part Two students are solving the same problem independently. If the probability that the first one solves the problem is $\frac{3}{5}$ and probability that the second solves the problem is $\frac{4}{5}$, what is the probability that at least one of them solves the problem? \\
    % https://csirhrdg.res.in/SiteContent/ManagedContent/ContentFiles/20190705160435567mathA_June2018.pdf
    % Que 18 part A
    \begin{oneparcheckboxes}
        \choice $\frac{17}{25}$.
        \choice $\frac{19}{25}$.
        \choice $\frac{21}{25}$.
        \correctchoice $\frac{23}{25}$.
    \end{oneparcheckboxes}
    %\vspace*{\stretch{1}}
    \part The number of heads, $X$, in one flip of a coin, is given by the following probability distribution.
    \begin{equation*}
        p(x)= (0.25)^x (0.75)^{1-x}, \quad x=0,1
    \end{equation*}
    $P\{X=1\}$ is \\
    \begin{oneparcheckboxes}
        \choice $0$.
        \correctchoice $0.25$.
        \choice $0.50$.
        \choice $0.75$.
    \end{oneparcheckboxes}
    %\vspace*{\stretch{1}}
    \part Let $X({\omega})=I_{A}(\omega)$ for some $A\in \mathcal{S}$. Where, $I_A$ is an indicator function of $A$. Then $E[X]$ is \\
    \begin{oneparcheckboxes}
        \choice $0$.
        \correctchoice $P(A)$.
        \choice $1-P(A)$.
        \choice $1$.
    \end{oneparcheckboxes}
    %\vspace*{\stretch{1}}
    \part Consider the distribution
    \begin{equation*}
        p(x) =\frac{3-x}{3}, \quad x=1,2
    \end{equation*}
    If $g(x)=3x+5$, then $E[g(X)]$ is equal to \\
    % Q3 page 73 http://www.stat.ucla.edu/~dinov/courses_students.dir/02/Fall/STAT100A.dir/STAT_100.dir/Stat100FinalReview.pdf
    \begin{oneparcheckboxes}
        \choice $\frac{21}{3}$.
        \choice $\frac{22}{3}$.
        \choice $\frac{23}{3}$.
        \correctchoice $\frac{24}{3}$.
    \end{oneparcheckboxes}
    %\vspace*{\stretch{1}}
    %\answerline
\end{parts}
}\label{que-2}
\newpage
\question[10]{Fill in the blanks. %\textbf{NO JUSTIFICATION IS REQUIRED.}
\begin{parts}
    \part Let $X$ be an RV with a Cauchy PDF
    \begin{equation*}
        f(x)=\frac{1}{\pi}\frac{1}{1+x^2}, \quad -\infty < x < \infty
    \end{equation*}
    The median of the RV $X$ is $x=$\fillin[0] %example 6 page 83 Rohatgi
    %\vspace*{\stretch{1}}
    \part Let 
    \begin{equation*}
        f(x) =\begin{cases}
            e^{-x}, & \text{if } x\geq 0 \\
            0, & \text{elsewhere}
        \end{cases}
    \end{equation*}
    Compute $E[X]$ if it exists; write "does not exist" otherwise.
    \begin{equation*}
        E[X]=\fillin[1]
    \end{equation*}
    %\vspace*{\stretch{1}}
    \part Let the random variable $X$ have moment generating function $M(t)=\frac{1}{(1-t)^2}$ for $t<1$. The third moment of $X$ about the origin is \fillin[24] 
    %tutorial 8 (4)
    %\vspace*{\stretch{1}}
    \part If $Z\sim N(0,1)$, what is the value of constant $c$ such that $P(|Z| \leq c)=0.95$?
    \begin{equation*}
        c=\fillin[1.96]
    \end{equation*}
    %tut 8 (9) PR
    %\vspace*{\stretch{1}}
    \part If $X \sim N(3,16)$ then $P\{4 \leq X \leq 8\}=\fillin[0.2957]$
    %\vspace*{\stretch{1}}
    \part In sampling from a non-normal distribution with a variance of $25$, the sample size must be \fillin[100] so that the length of the $95\%$ confidence interval for the mean will be $1.96$.
    %\vspace*{\stretch{1}}
    \part If $P(A)=p_1$, $p(B)=p_2$, and $P(A \cap B)=p_3$, then $P(A^c \cup B)$ is \fillin[$1 - p_1 + p_3$]
    %If $X \sim N(7,4)$, then $P\left\{15.364\leq (X-7)^2 \leq 20.095\right\}=$\fillin
    %\vspace*{\stretch{1}}
    \part Let $X$ be defined on $\left(\Omega, \mathcal{S}, P\right)$ by
    \begin{equation*}
        X\left(\omega\right) = c, \quad \text{for all } \omega \in \Omega 
    \end{equation*}
    Then $F(x)=\fillin$
    \begin{equation*}
        F(x)=
        \begin{cases}
            & \\
            &  \\ 
            & 
        \end{cases}
    \end{equation*}
    (Here, $c$ is a constant, and $F$ is DF; write cases for the $F$ in the above space. )
    %Example 1 Page 47 Rohatgi
    %\vspace*{\stretch{1}}
    \part If $X \sim B(n,p)$ the parameter space for $p$ is \fillin[$0<p<1$]
    %\vspace*{\stretch{1}}
    %Write the parameter space for the $\theta$ if
    %\begin{subparts}
        %\subpart If $X \sim EXP(\theta)$ the parameter space for $\theta$ is \fillin.
        %\newline
        %\subpart .
    %\end{subparts}
    \part Let $X$ have the uniform distribution on the first $N$ natural numbers; that is, let 
    \begin{equation*}
        P\left\{X=k\right\} =\frac{1}{N}, \quad k=1,2,\ldots, N
    \end{equation*}
    Then %example 2 Rohatgi page 73
    \begin{equation*}
        E(X^2)= \fillin
    \end{equation*}
    %\vspace*{\stretch{1}}
    %\begin{subparts}
        %\subpart $E(X)=$ \fillin
       % \subpart $E(X^2)=$ \fillin
    %\end{subparts}
\end{parts}
    } \label{que-3}
\newpage
\question{Answer the following questions.
\begin{parts}
    \part[2] Two players, $A$ and $B$ play a coin-tossing game. $A$ gives $B$ one dollar if a head turns up; otherwise, $B$ pays $A$ one dollar. If the probability that the coin shows a head is $p$, find the expected gain of $A$.
    \begin{solution}
        %Rohatgi page 73 Rohatgi
        $X$: the gain of $A$
        \begin{equation*}
            P\{X=1\}=P\{\text{tails}\}=1-p, \quad P\{X=-1\}=p.
        \end{equation*}
        and
        \begin{equation*}
            E[X]=1-p-p=1-2p.
        \end{equation*}
    \end{solution}
    %\makeemptybox{3.5in}
    \part[2] Prove that if events $A$ and $B$ are independent, then events $A$ and $B^c$ (read as $B$ complement) are also independent.
    \begin{solution}
        \begin{equation*}
            A = (A \cap B^c) \cup (A \cap B).
        \end{equation*}
        Thus,
        \begin{align*}
            P(A) &= P\left((A \cap B^c) \cup (A \cap B)\right) \\
            &= P((A \cap B^c))+ P((A \cap B)) \\
            \implies  P((A \cap B^c)) &= P(A)- P((A \cap B)) \\
            &= P(A)-P(A)P(B) \\
            &= P(A)(1-P(B)) \\
            &= P(A)P(B^c).
        \end{align*}
    \end{solution}
    \part[2] A normal population has a mean of $0.1$ and standard deviation of $2.1$. Find the probability that mean of a sample of size $900$ will be negative.
    \begin{solution}
        \begin{align*}
            Z&=\frac{\bar{x}-\mu}{\frac{\sigma}{\sqrt{n}}} \\
            &= \frac{\bar{x}-0.1}{\frac{2.1}{30}} \\
            \implies \bar{x} &= 0.1+0.07Z, \quad Z \sim N(0,1).
        \end{align*}
        Now,
        \begin{align*}
            P\{\bar{x}<0\} &= P\left\{0.1+0.07Z<0\right\}\\
            &= P\left\{Z < \frac{-0.10}{0.07}\right\} \\
            P\left\{Z < -1.43\right\} %\\
            &=0.0764
        \end{align*}
        %Tutorial 11. Que. (8)
        %%12.20 https://www.dcpehvpm.org/E-Content/Stat/FUNDAMENTAL%20OF%20MATHEMATICAL%20STATISTICS-S%20C%20GUPTA%20&%20V%20K%20KAPOOR.pdf
    \end{solution}
    \part[2] Let $X$ be an RV with PDF
    \begin{equation*}
        f(x) =\begin{cases}
            1 & \text{if } 0 < x < 1 \\
            0 & \text{otherwise.}
        \end{cases}
    \end{equation*}
    Let $Y=-2\log{X}$. Find the PDF of $Y$.
    \begin{solution}
        If $y=-2\log{x}$, then $x=e^{\frac{y}{2}}$, and 
        \begin{align*}
            h(y) &= \left|\frac{1}{2}e^{-\frac{y}{2}}\right| \cdot 1, & 0< e^{\frac{-y}{2}}<1,\\
            &= \begin{cases}
                \frac{1}{2}e^{-\frac{y}{2}} & 0<y<\infty \\
                0 & \text{otherwise.}
            \end{cases}
        \end{align*}
    \end{solution}
    \part[2] Compute the moment generating function for a standard Normal random variable.
    % https://www.le.ac.uk/users/dsgp1/COURSES/MATHSTAT/6normgf.pdf
    \begin{solution}
        \begin{align*}
            \phi(t)=E(e^{zt})&=\int e^{zt}\frac{1}{\sqrt{2 \pi}}e^{-\frac{1}{2}z^2}dz \\
            &= e^{-\frac{-t^2}{2}}.
        \end{align*}
    \end{solution}
    
    %\part[2] 
\end{parts}
} \label{que-4}
%\vspace*{\stretch{1}}
\question{Answer the following questions.
\begin{parts}
    \part[3] Show that if $X$ is random variable, $X^2$ is also  random variable.
    \part[3] Let $X$ be an $RV$ with $PDF$
        \begin{equation*}
            f(x) = \begin{cases}
            0 & \text{if } x \leq 0, \\
            \frac{1}{2} & \text{if } 0< x \leq 1,\\
             \frac{1}{2x^2} & \text{if } 1 < x < \infty.
                    \end{cases} 
        \end{equation*}
    Find the $PDF$ of $RV$ $\frac{1}{X}$.
    \begin{solution}
        $X$ and $Y=\frac{1}{X}$ are identically distributed.
        \begin{equation*}
            h(y) = \begin{cases}
                0 & \text{if } y \leq 0, \\
            \frac{1}{2} & \text{if } 0< y \leq 1,\\
             \frac{1}{2y^2} & \text{if } 1 < y < \infty.
            \end{cases}
        \end{equation*}
    \end{solution}
    \part[3] If $X$ is uniform on the interval from $0$ to $3$, what is the probability that the quadratic equation $4t^2 + 4tX + X + 2 = 0$ has real solutions?
    \begin{solution}
        %Tutorial 10. Que 11.
        %%t8 16
        $X\sim U(0,3)$. PDF of $X$ is
        \begin{equation*}
            f(x) = \begin{cases}
                \frac{1}{3} & 0 \leq x \leq 3 \\
                0 & \text{otherwise.}
            \end{cases}
        \end{equation*}
        The discriminant $16X^2-16(X+2)=(X-2)(X+1) \geq 0$.
        \begin{align*}
            P((X-2)(X+1) \geq 0) &= P(X \leq -1 \text{ or }X \geq 2) \\
            &= P(X \leq -1)+P(X \geq 2) \\
            &=\frac{1}{3}
        \end{align*}
    \end{solution}
    %\part[3] 
    \part[3] If $X_1, X_2, \ldots, X_n$ is a random sample from a distribution with density function
    \begin{equation*}
        f(x; \theta) = \begin{cases}
            (1-\theta)x^{-\theta}, & \text{if } 0<x<1 \\
            0, & \text{elsewhere.}
        \end{cases}
    \end{equation*}
    what is the maximum likelihood estimator of $\theta$?
    \begin{solution}
        %Lecture 28-N Example 3
        \begin{align*}
            L(\theta) &= \prod_{i=1}^{n}f(x_i, \theta)\\
           \text{Therefore, } \ln{L(\theta)} &= n \ln{(1-\theta)}-\theta \sum_{i=1}^n \ln{x_i}.
        \end{align*}
        \begin{equation*}
            \frac{d \ln{L(\theta)}}{d \theta} = \frac{-n}{1-\theta}-\sum_{i=1}^{n}\ln{x_i}
        \end{equation*}
        At $\theta = 1+ \frac{1}{\bar{\ln{x}}}$, $\frac{d \ln{L(\theta)}}{d \theta}=0$. This $\theta$ can be shown to be maximum by the second derivative test
    \end{solution}
    \part[3] Let $X_1, X_2, \ldots, X_n$ be a random sample of size $11$ from a normal distribution with unknown mean $\mu$ and variance $\sigma^2=9.9$. If $\sum_{i=1}^{11}x_i=132$, then what is the $95\%$ confidence interval for $\mu$?
    \begin{solution}
        %lecture-30-31-N example 9
        Since, $X_i \sim N(\mu, 9.9)$, the confidence interval for $\mu$ is given by
        \begin{equation*}
            \left[\bar{x}-\left(\frac{\sigma}{\sqrt{n}}\right)z_{\frac{\alpha}{2}},\bar{x}+\left(\frac{\sigma}{\sqrt{n}}\right)z_{\frac{\alpha}{2}}\right]
        \end{equation*}
        We have, $\bar{x}=12$, $\sqrt{\frac{\sigma^2}{n}}=\sqrt{0.9}$, $\alpha=0.05$, and $z_{\frac{\alpha}{2}}=1.96$. Using this information in the expression of the confidence interval for $\mu$, we get
        \begin{equation*}
            \left[10.141, 13.859\right]
        \end{equation*}
    \end{solution}
    %\part[3] Show that if $X$ is random variable, $X^2$ is also  random variable.
\end{parts}
} \label{que-5}
%\vspace*{\stretch{1}}
\question[5]{A dice is thrown $9000$ times and throw of $3$ or $4$ is observed $3240$ times. Show that dice cannot be regarded as unbiased one and find the limits between which probability of a throw of $3$ or $4$ lies.
    \begin{solution}
        %12.1 https://www.dcpehvpm.org/E-Content/Stat/FUNDAMENTAL%20OF%20MATHEMATICAL%20STATISTICS-S%20C%20GUPTA%20&%20V%20K%20KAPOOR.pdf
        If the coming of $3$ or $4$ is called a success, then in usual notations we are given,
        \begin{equation*}
            n=9000, \quad X= \text{Number of Success} = 3240
        \end{equation*}
        Under the null hypothesis $(H_0)$ that the dice is an unbiased one. we get
        \begin{equation*}
            P=\text{Probability of success}=\text{Probability of getting a 3 ot 4} =\frac{1}{3}
        \end{equation*}
        Alternate hypothesis: $H_1: p \neq \frac{1}{3}$. (i.e. dice is biased.)
        \begin{equation*}
            \text{We have } Z= \frac{X-nP}{\sqrt{nQP}} \sim N(0,1)
        \end{equation*}
        Now, $Z=5.36$. Since $|Z|>3$, $H_0$ is rejected, we conclude that dice is almost certainly biased.
        \par Since dice is not unbiased, $p\neq \frac{1}{3}$. The probable limits for $P$ are given by.
    \end{solution}} \label{que-6}
\end{questions}
\end{document}
